This handout shows how the Envelope theorem is used to derive the
consumption Euler equation in a multiperiod optimization problem with
geometric discounting and intertemporally separable utility.

The consumer's goal from the perspective of date $\tNow$ is to \opt{MarginNotes}{\marginpar{\tiny
Define `geometric discounting.'
}}
\begin{equation}
        \max \sum_{n=0}^{\TEnd-\tNow} \Discount^{n} \uFunc(c_{\tNow+n})
\end{equation}
subject to the dynamic budget constraint
\begin{equation}\begin{gathered}\begin{aligned}
 {m}_{\tNow+1} & =  ({m}_{\tNow}-{c}_{\tNow})\Rfree + {y}_{\tNow+1}. \label{eq:xaccum}
\end{aligned}\end{gathered}\end{equation}

The problem can be written in Bellman equation form as
\begin{equation}\begin{gathered}\begin{aligned}
        \vFunc_{\tNow}({m}_{\tNow}) & =  \max_{\{{c}_{\tNow}\}} ~\uFunc({c}_{\tNow}) + \Discount \vFunc_{\tNow+1}(({m}_{\tNow}-{c}_{\tNow})\Rfree+{y}_{\tNow+1}). \label{eq:bellman}
\end{aligned}\end{gathered}\end{equation}

The first order condition for \eqref{eq:bellman} can be written
as
\begin{equation}\begin{gathered}\begin{aligned}
        0 & =
        \uFunc^{\prime}({c}_{\tNow})+\overbrace{\left(\frac{d {m}_{\tNow+1}}{d {c}_{\tNow}}\right)}^{=-\Rfree~\mbox{from \eqref{eq:xaccum}}}
        \Discount \vFunc_{\tNow+1}^{\prime}({m}_{\tNow+1})
\\  \uFunc^{\prime}({c}_{\tNow}) & =  \Rfree \Discount \vFunc_{\tNow+1}^{\prime}({m}_{\tNow+1}), \label{eq:cfoc}
\end{aligned}\end{gathered}\end{equation}
and we can define a function $\cFunc_{\tNow}({m})$ that returns the
${c}_{\tNow}$ that solves the max problem for any given ${m}_{\tNow}$.  That is,
for ${c}_{\tNow}=\cFunc_{\tNow}({m}_{\tNow})$ the first order condition \eqref{eq:cfoc}
will hold so that
\begin{equation}\begin{gathered}\begin{aligned}
        \uFunc^{\prime}(\cFunc_{\tNow}({m}_{\tNow})) - \Rfree \Discount \vFunc_{\tNow+1}^{\prime}(({m}_{\tNow}-\cFunc_{\tNow}({m}_{\tNow}))\Rfree+{y}_{\tNow+1}) & =  0. \label{eq:cstarfoc}
\end{aligned}\end{gathered}\end{equation}

Now define a function \opt{MarginNotes}{\marginpar{\tiny
The _ indicates it's a (weak) lower bound
}}
\begin{equation}\begin{gathered}\begin{aligned}
        \underline{\vFunc}_{\tNow}({m}_{\tNow},{c}_{\tNow}) & =  \uFunc({c}_{\tNow})+\Discount \vFunc_{\tNow+1}(({m}_{\tNow}-{c}_{\tNow})\Rfree+{y}_{\tNow+1})
\end{aligned}\end{gathered}\end{equation}
with partial derivatives
\begin{equation}\begin{gathered}\begin{aligned}
        \underline{\vFunc}^{c}_{\tNow}({m}_{\tNow},{c}_{\tNow}) \equiv \left(\frac{\partial \underline{\vFunc}_{\tNow}}{\partial {c}_{\tNow}}\right) & =
        \uFunc^{\prime}({c}_{\tNow}) - \Rfree\Discount \vFunc_{\tNow+1}^{m}(({m}_{\tNow}-{c}_{\tNow})\Rfree+{y}_{\tNow+1}) \label{eq:nuc}
\end{aligned}\end{gathered}\end{equation}
\begin{equation}\begin{gathered}\begin{aligned}
      \underline{\vFunc}^{{m}}_{\tNow}({m}_{\tNow},{c}_{\tNow}) \equiv \left(\frac{\partial \underline{\vFunc}_{\tNow}}{\partial {m}_{\tNow}}\right) & =  \Rfree \Discount \vFunc_{\tNow+1}^{{m}}({m}_{\tNow+1}) \label{eq:nuprimevtp1}
\end{aligned}\end{gathered}\end{equation}
and note that by definition
\begin{equation}\begin{gathered}\begin{aligned}
        \vFunc_{\tNow}({m}_{\tNow}) & =  \underline{\vFunc}_{\tNow}({m}_{\tNow},\cFunc_{\tNow}({m}_{\tNow})).
\end{aligned}\end{gathered}\end{equation}

The Chain Rule of differentiation tells us that
\begin{equation}\begin{gathered}\begin{aligned}
        \vFunc^{\prime}_{\tNow}({m}_{\tNow})\equiv \vFunc_{\tNow}^{{m}}({m}_{\tNow}) \equiv \left(\frac{d \vFunc_{\tNow}}{d {m}_{\tNow}}\right) & =  \underline{\vFunc}_{\tNow}^{{m}}({m}_{\tNow},\cFunc_{\tNow}({m}_{\tNow})) + \left(\frac{\partial \cFunc_{\tNow}({m}_{\tNow})}{\partial {m}_{\tNow}}\right)\underline{\vFunc}_{\tNow}^{c}({m}_{\tNow},\cFunc_{\tNow}({m}_{\tNow})). \nonumber
\\ \label{eq:partdiff}
\end{aligned}\end{gathered}\end{equation}

Here's the key insight:  The assumption that consumers are optimizing
means that we will always be evaluating the value function and its
derivatives at a ${c}_{\tNow}$
that satisfies the first-order optimality condition \eqref{eq:cstarfoc}.\footnote{Unless there is some constraint that prevents the consumer from choosing this optimum - like a liquidity constraint.}  Thus
we have from \eqref{eq:nuc} that
\begin{equation}\begin{gathered}\begin{aligned}
        \underline{\vFunc}_{\tNow}^{c}({m}_{\tNow},\cFunc_{\tNow}({m}_{\tNow})) & =  \uFunc^{\prime}(\cFunc_{\tNow}({m}_{\tNow})) - \Rfree\Discount \vFunc^{\prime}_{\tNow+1}(({m}_{\tNow}-\cFunc_{\tNow}({m}_{\tNow}))\Rfree+{y}_{\tNow+1})
\\ & =  0.
\end{aligned}\end{gathered}\end{equation}
This means that the second term in \eqref{eq:partdiff}
is always equal to zero, so from \eqref{eq:nuprimevtp1} we obtain
\begin{equation}\begin{gathered}\begin{aligned}
        \vFunc^{\prime}_{\tNow}({m}_{\tNow}) & =  \Rfree \Discount \vFunc_{\tNow+1}^{\prime}({m}_{\tNow+1}). \label{eq:eulerv}
\end{aligned}\end{gathered}\end{equation}

Now notice that the RHS's of \eqref{eq:cfoc} and \eqref{eq:eulerv} are
identical, so we can equate the left hand sides,
\begin{equation}\begin{gathered}\begin{aligned}
        \vFunc_{\tNow}^{\prime}({m}_{\tNow}) & =  \uFunc^{\prime}({c}_{\tNow})
\end{aligned}\end{gathered}\end{equation}
and since a corresponding equation will hold in period $t+1$
we can rewrite \eqref{eq:eulerv} as
\begin{equation}\begin{gathered}\begin{aligned}
        \uFunc^{\prime}({c}_{\tNow}) & =  \Rfree \Discount \uFunc^{\prime}({c}_{\tNow+1}).
\end{aligned}\end{gathered}\end{equation}

The general principle can be condensed into a rule of thumb by
realizing that the Envelope theorem will always imply that the total
derivative of a value function with respect to any choice variable
must be equal to zero for optimizing consumers (because the first
order condition holds).  Thus we could have obtained the result
immediately by treating ${c}_{\tNow}$ as though it were a constant
(that is, treating the problem as though
$\cFunc_{\tNow}^{\prime}(m_{\tNow})=0$) and taking the derivative of
Bellman's equation with respect to ${m}_{\tNow}$ directly.  This leads
immediately to the key result:
\begin{equation}\begin{gathered}\begin{aligned}
        \vFunc_{\tNow}({m}_{\tNow}) & =  \uFunc(\cFunc({m}_{\tNow})) + \Discount \vFunc_{\tNow+1}(({m}_{\tNow}-\cFunc_{\tNow}({m}_{\tNow}))\Rfree+{y}_{\tNow+1})
\\      \vFunc_{\tNow}^{\prime}({m}_{\tNow}) & =  \Discount \Rfree \vFunc_{\tNow+1}^{\prime}({m}_{\tNow+1}).
\end{aligned}\end{gathered}\end{equation}

\opt{MarginNotes}{\marginpar{\tiny Using the figure, discuss the
    consequences of an increase in m, and show that the increase in
    utility is about the same whether you consume all or none of it.}}

\begin{figure}
\caption{Illustration of the Envelope Theorem at Alternative Values of $\mRat$}
\includegraphics[width=6in]{../Figures/Envelope}
\end{figure}

